\documentclass{article}
\usepackage{tikz}
\usetikzlibrary{angles, quotes}
\usepackage{amssymb}
\usepackage{pgfplots}
\pgfplotsset{width=10cm,compat=1.9}
\usepgfplotslibrary{statistics}
\usepackage[left=2.54cm,right=2.54cm,top=2.54cm,bottom=2.54cm]{geometry}
\usepackage{tasks} %Paquete necesario para la producción de items horizontales para algunos ejercicios de tarea empleando el comando /begin{multicols}}{#}.../end{multicols} para ayudarnos a reducir las páginas
\usepackage{pdflscape} %comando que permite cambiar a orientación horizontal las hojas%
\usepackage{soul} %Paquete para permitir la compilación de acentos usando el comando {\hl{}}}\\
\sethlcolor{yellow(munsell)} %comando para definir el color para el subrayado usando el comando \hl
\usepackage[utf8]{inputenc}
\usepackage[spanish]{babel}
\usepackage{icomma}
\usepackage{ragged2e}
\usepackage{siunitx}
\usepackage{url}
\usepackage[colorlinks=true, urlcolor=blue,  linkcolor=black, citecolor=green]{hyperref}
\usepackage{pdfpages}
\usepackage{blindtext} %Paquete que genera texto ficticio (dummy text) con el comando: \blindtext
\usepackage{cancel} %in the preamble gives you four different modes of striking through: \cancel{text to cancel} draws a diagonal line (slash) through its argument, \bcancel{text to cancel} uses the negative slope (a backslash), \xcancel{text to cancel} draws an X (actually \cancel plus \bcancel), \cancelto{〈value〉}{〈expression〉} draws a diagonal arrow through the 〈expression〉pointing to the 〈value〉 (math-mode only)
\usepackage{csquotes}
\usepackage{afterpage}
\usepackage{parskip} 
\usepackage{float}
\usepackage{enumitem}
\usepackage{multicol}%Paquete que permite la creación aislada de columnas de texto. De esta forma se reduce la cantidad de páginas en nuestro documento
\usepackage{lipsum} %Paquete que genera texto de relleno al igual que \usepackage{blindtext}, se usa con el comando \lipsum[#-#] los #(hashtags) delimitan el número de páginas que deseamos ocupar del paquete lipsum para usarlos como relleno. 

% Definir el comando personalizado para la nota
\newcommand{\mynote}[1]{%
\begin{tikzpicture}[baseline=-0.75ex]
\node[draw, circle, fill=Apple Green!60, inner sep=2pt] (note) {\textbf{Nota:}};
\end{tikzpicture}%
\ \textit{#1}%
}

\newenvironment{Figura}
{\par\medskip\noindent\minipage{\linewidth}}
{\endminipage\par\medskip}
\usepackage{caption}
\usepackage[
backend=biber,
sorting=none,
url=true
]{biblatex} %bibliografía
\addbibresource{biblio.bib}
% Establecer el espacio entre las entradas de la bibliografía
\setlength{\bibitemsep}{1\baselineskip} % Puedes ajustar el valor según tus preferencias
\usepackage{amsmath,amsthm,amssymb,amsfonts}
\usepackage{pifont} %Permite la compilación del comando \xmark: es el símbolo de la tachita
\usepackage{mathtools} %permite la compilación de símbolos para matrices
\usepackage{empheq} %Paquete que se relaciona con \usepackage[most]{tcolorbox} para la creación de cuadros/cajas de colores para delimitar los resultados matemáticos o líneas de texto usando la línea de comando: \begin{empheq}[box={\mymath[colback=orange(webcolor)!70,drop lifted shadow]}]{equation*} \end{empheq}
\usepackage[most]{tcolorbox} %Paquete que se relaciona con \usepackage{empheq} para la creación de cuadros/cajas de colores para delimitar los resultados matemáticos o líneas de texto
\renewcommand{\qedsymbol}{$\blacksquare$}
\usetikzlibrary{positioning,decorations.pathreplacing} %Paquete que permite la compilación de llaves para cuadros sinópticos (brace diagram) 
\usepackage{schemata} %The schemata package is designed just for these brace diagrams
\usepackage{booktabs}

\newcommand\AB[2]{\schema{\schemabox{#1}}{\schemabox{#2}}} %comando o paquete necesario para crear cuadros sinópticos

\newtcbox{\mymath}[1][]{%
nobeforeafter, math upper, tcbox raise base,
enhanced, colframe=blue!30!black,
colback=blue!30, boxrule=1pt,
#1} %Comando importante para encasillar los resultados matemáticos en cajas de diferentes colores, se relaciona con los paquetes \usepackage{empheq} y \usepackage[most]{tcolorbox} 

\newenvironment{sysmatrix}[1]
{\left(\begin{array}{@{}#1@{}}}
{\end{array}\right)}
\newcommand{\ro}[1]{%
\xrightarrow{\mathmakebox[\rowidth]{#1}}%
}
\newlength{\rowidth}% row operation width
\AtBeginDocument{\setlength{\rowidth}{3em}}  %Comando importante para laproducción de lineas de operaciones entre matrices, método de Gauss o Gauss Jordan se relaciona con \newenvironment{sysmatrix}

%formato para cambiar el horario a español
\usepackage[spanish]{datetime2}
\DTMsetdatestyle{spanish}

\renewcommand{\today}{\DTMdisplaydate{\the\year}{\the\month}{\the\day }{-1}}
%%%%%%%%%%%%%%%%%%%%%%%%%%%%%%%%%%%%%%%%%

\usepackage{datetime}
\newcommand{\mycurrenttime}{\xxivtime}

%%%%%%%%%%%%%%%%%%%%%%%%%%%%%%%%%%%%%%%%%%%%%%%%%%%%%%%%%%%%%%%%%%%%%%%%%%%%%%%%%%%%%%%%%%%%%%%%%%%%%%%%%%%%%%%%%%%%%%%%%%%%%%%%%%%%%%%%%%%%%%%%%


%%%%%%%%%%%%%Caja de color para el título%%%%%%%%%%%%%%%%%%%%%%%%%%%%%%%%%%%%%%%%%%%%%%%%%%%%%%

\definecolor{myframecolor}{RGB}{1, 45, 75} % Prussian Blue
\definecolor{myboxcolor}{RGB}{0, 107, 92} % Apple Green

\newtcolorbox{mybox}{
enhanced,
colback=myboxcolor!7, % Color de fondo del cuadro
colframe=myframecolor, % Color del marco del cuadro
arc=0pt,
boxrule=1pt,
borderline west={2mm}{-10mm}{myframecolor}, % Borde en el lado izquierdo
sharp corners=southwest,
width=\linewidth,
before=\par\vspace{\bigskipamount}, % Espacio antes del cuadro
after=\par\vspace{\bigskipamount} % Espacio después del cuadro
}


\newcommand{\euler}{e} %Comando para producir letra e de euler
\newenvironment{remark}{\par\vfill\footnotesize % Vertical white space above the remark and smaller font size
\begin{list}{}{
	\leftmargin=80pt % Indentation on the left
	\rightmargin=60pt}\item\ignorespaces % Indentation on the right
\makebox[-2.5pt]{\begin{tikzpicture}[overlay]
		\node[draw=Horizon!60,line width=2.5pt,circle,fill=Horizon!25,font=\sffamily\bfseries,inner sep=4pt,outer sep=0pt] at (-19pt,5pt){\textcolor{Horizon}{Nota}};\end{tikzpicture}} % Blue Nota in a circle
\advance\baselineskip -1pt}{\end{list}\vskip5pt} % Tighter line spacing and white space after remark
\usepackage{graphicx}

\usepackage{titling}

\input{colores.tex}

%%Producir signos de cita textual``Asignatura''

\newcommand{\xmark}{\ding{55}} %%%comando que genera la tachita

\newenvironment{MyColorPar}[1]{%
\leavevmode\color{#1}\ignorespaces%
}{%
}%

%%%%%%%%%%%%%%%%%%% Título de la tarea, Nombre de alumno

\title{Clase \# 43 - Mecánica Teórica}
\author{\emph{JA} $\boldsymbol{\mid}$}
\date{\today}

%%%%%%%%%%%%%%%%%%% Datos de la Materia

\usepackage{fancyhdr}
\fancypagestyle{plain}{%  the preset of fancyhdr 
\fancyhf{} % clear all header and footer fields
\fancyfoot[R]{
\includegraphics[width=2cm]{LogoFCFMBUAP (1).png}}
\fancyfoot[L]{{\bfseries{\thedate{}}} a las {\bfseries{\mycurrenttime{}}} horas{} (GMT-6, H. Puebla de Zaragoza, Pue)}
\fancyhead[L]{Mecánica teórica}
\fancyhead[R]{\theauthor}
}
\makeatletter
\def\@maketitle{%
\newpage
\null
\vskip 1em%
\begin{center}%
\let \footnote \thanks
{\LARGE \@title \par}%
\vskip 1em%
%{\large \@date}%
\end{center}%
\par
\vskip 1em}
\makeatother

\usepackage{lipsum}  
\usepackage{cmbright}


%%%%%%%%%%%%%%%%%%%%%%%%%%%%%%%%%%%%%%%%%%%%%%%%%%%%%%%%%%%%%%%%%%%%%%%%%%%%%%%%%%%%%%%%%%%%%%%%%%%%%%%%%%%%%%%%%%%%%%%%%%%%%%%%%%%%%%%%%%%%%%%%%%%%%%%%%%%%%%%%%%%%%%%%%%%%%%%%%%%%%%%%%%%%%%%%%%%%%%

\begin{document}



\maketitle


%%%%%%%%%%%%%%%%%%% Datos del profesor, campus y aula

\begin{mybox}
\noindent\begin{tabular}{@{}ll}
	{\bfseries{Profesor}} & Dr. Gilberto Silva Ortigoza\\
	\textcolor{prussianblue}{\bfseries{Campus}} / \textcolor{trueblue}{\bfseries{Facultad}}  & \textcolor{prussianblue}{\bfseries{C.U. BUAP}} /  \textcolor{trueblue}{\bfseries{FCFM}} \\
	\textcolor{prussianblue}{\bfseries{Edificio}} /  \textcolor{trueblue}{\bfseries{EMA4}}    & 	\textcolor{prussianblue}{\bfseries{FM9}} /  \textcolor{trueblue}{\bfseries{302}}   
\end{tabular} 
\end{mybox} \vspace{1cm}

\tableofcontents

\vspace{2cm}

\section{Clase del día de hoy}


Tenemos:


El moemnto angular $\vec{L} \equiv \vec{r} \times \vec{p} \quad \text{es el generador}$ de rotaciones espaciales.

\begin{center}
\includegraphics[width=0.85\linewidth]{1.png} % <-- Sustituye esta ruta
\\[2mm]
\textbf{Figure 1}
\end{center}	


Vamos a probar que si tenemos:


\[  Q_{i} = q_{i} + \epsilon \frac{ \partial G  }{\partial p_{i}}, \quad G = \vec{L} \cdot z \]



Tenemos:


\[  P_{i} = p_{i} - \epsilon \frac{ \partial G  }{\partial q_{i}} \]


Tenemos

\[ \vec{r}(t) = x(t)\hat{x} + y(t) \hat{y} + z(t)\hat{z} \]


Tenemos:


\[  q_{1} = x; \quad q_{2} = y, \quad q_{3} = z \]


\[ \vec{A} = A_{x} \hat{x} + A_{y} \hat{y} + A_{z} \hat{z} \]

\emph{Las rotaciones vienen determinados por 3 parámetros.} Un cuerpo rígido experimenta una rotación y traslación.


El movimiento más general que va a experimentar un cuerpo rígido es un movimiento de \underline{traslación} y \underline{rotación}.


Tenemos:



El número de grados de libertad que tiene un cuerpo rígido en el espacio es:


\[ n = 3N -\mathcal{K}\]


Anre una rotación en el eje en el eje $x-y$, la componente $z$ no se ve afectada.


Tenemos:


\[ \bar{A}_{x} = A_{x} \cos \phi + A_{y} \sen \phi  \]

Tenemos:


\[ \bar{A}_{y} = -A_{x}\sen\phi + A_{y}\cos\phi  \]


\[\bar{A}_{z} = {A}_{z} \]


Entonces ante una rotación las componente se transforman (no infinitesimal)

Tenemos:


\[  ( \bar{A}_{x}, \bar{A}_{y}, A_{z} ) \]


Entonces tenemos (\underline{rotación infitesimal, hacemos una expansión en serie alrededor del $0$ para $\phi$}).


\[ f(x) \approx \sum\limits_{i=1}^{n} \frac{  f^{(n) }a }{ n!  } (x-a)^{n}    \]


Tomando \[ a=0, \quad \text{cuando hacemos eso se llama \textbf{serie de Mclaurin}} \]


Tenemos:


\textbf{A primer orden} (dado que las transformaciones infinitesimales vamos a trabajr en primer orden en $\epsilon$)


Tenemos:

 \[  \cos \phi = \cos 0 + \left. \frac{d \cos \phi }{ d \phi} \right|_{\phi=0} \phi + \ldots     \]

Tenemos:

\[ \cos \phi = 1 - \left. \sen \phi \right|_{\phi=0} \phi\]

\[ \approx 1 \]



Tenemos:


\[   \sen \phi = \sen 0 + \left. \frac{d \sen \phi }{ d \phi} \right|_{\phi=0} \phi + \ldots   \]


\[ 0 + \left.\cos \phi \right|_{\phi=0} \phi \]

Tenemos:


\[\boxed{ \equiv \phi} \]



\begin{center}
\includegraphics[width=1\linewidth]{2.png} % <-- Sustituye esta ruta
\\[2mm]
\textbf{Figure 2}
\end{center}	

Tenemos:

\begin{center}
\includegraphics[width=1\linewidth]{3.png} % <-- Sustituye esta ruta
\\[2mm]
\textbf{Figure 3}
\end{center}	

Hallando las otras 3 ecuaciones:



\begin{center}
\includegraphics[width=1\linewidth]{4.png} % <-- Sustituye esta ruta
\\[2mm]
\textbf{Figure 4}
\end{center}	


\textbf{Obtención de $G$}


SI realizamos bien las cuentas, tendremos nosotros. 

\[  G = q_{2}p_{1} - q_{1}p_{2}\]

\[ = yp_{x}-xp_{y} = (\vec{r} \times \vec{p})_{z}  \]


\[ = \vec{L} \cdot \hat{z} = L_{z} \]


\begin{center}
\includegraphics[width=1\linewidth]{5.png} % <-- Sustituye esta ruta
\\[2mm]
\textbf{Figure 5}
\end{center}	

\subsection{Importante}


SI denotamos:


\[ L_{i} = (L_{x},L_{y},L_{z}) \]


Tenemos:


\[ \vec{L} \equiv \vec{r} \times \vec{p}\]


\[  \left[ L_{i}, L_{j}   \right] = \epsilon_{ijk}L_{k} \]

\[  \boxed{ \left[ L_{i}, V_{j}    \right] = i \hbar \epsilon_{ijk} V_{k}}  \]



\[ V_{i} = (v_{1}, v_{2}, v_{3}) \]

\[ = (v_{x}, v_{y}, v_{z}) \]

\textbf{Mañana demuestra el resultado anterior el profe Ortigoza.}



\end{document}